\chapter{Retrieval: Dense Search, Reranking and Decomposition}
\label{ch:retrieval}

Retrieval decides which five passages the language model is allowed to read. If the right passage is not among them, no prompt, grader or policy can produce a correct cited answer. This chapter covers the three retrieval techniques in the code: dense search, which every profile uses; cross-encoder reranking, which ships; and query decomposition with reciprocal rank fusion, which was built, measured and rejected.

\section{The retrieval stage as a whole}
\label{sec:retrieve-stage}

All retrieval runs through one method, \code{Pipeline.retrieve}, and every optional step is switched by a profile flag.

\begin{lstlisting}[language=Python]
def retrieve(self, question: str, trace: QueryTrace) -> list[RetrievedChunk]:
    sub_questions: list[str] = []
    if self.config.retrieval.use_decomposition:
        with trace.stage("decompose"):
            split = decompose(self.llm, question, model=self.config.generation_model)
        sub_questions = split.sub_questions
        trace.sub_questions = list(sub_questions)

    with trace.stage("retrieve"):
        if sub_questions:
            candidates = retrieve_decomposed(self.client, sub_questions, ...,
                                             limit=self.config.retrieval.top_k_retrieve,
                                             )[: self.config.retrieval.top_k_retrieve]
        else:
            candidates = retrieve_dense(self.client, question, ...,
                                        limit=self.config.retrieval.top_k_retrieve)
    trace.retrieved_chunk_ids = [c.chunk_id for c in candidates]      # the pool

    if self.config.retrieval.use_reranker and candidates:
        with trace.stage("rerank"):
            reordered = rerank(question, candidates,
                               model_name=self.config.retrieval.reranker_model)
        trace.rerank = rank_movement(candidates, reordered)
        candidates = reordered

    trace.ranked_chunk_ids = [c.chunk_id for c in candidates]         # the ordering
    context = candidates[: self.config.retrieval.top_k_context]         # the prompt

    if self.config.guardrails.check_retrieved_context and context:
        with trace.stage("guard_context"):
            screened = screen_context(context)
        trace.guard_context = screened.as_dict()
        context = screened.kept
    return context
\end{lstlisting}
\vspace{-4pt}{\raggedleft\ev{src/vidyarag/pipeline.py:84-142} (arguments abridged)\par}

\begin{figure}[htbp]
\centering
\resizebox{\textwidth}{!}{%
\begin{tikzpicture}[
  font=\sffamily\footnotesize,
  node distance=5mm and 7mm,
  proc/.style={draw, rounded corners=2pt, fill=blue!6, align=center, text width=30mm, minimum height=12mm},
  opt/.style={draw, dashed, rounded corners=2pt, fill=orange!8, align=center, text width=30mm, minimum height=12mm},
  rec/.style={draw, fill=yellow!14, align=center, text width=30mm, font=\sffamily\scriptsize},
  arr/.style={-Stealth, thick}
]
\node[opt] (dec) {\textbf{decompose}\\(if \code{use_decomposition})\\Gemini, JSON};
\node[proc, right=of dec] (ret) {\textbf{retrieve}\\dense top 20, or one top 20 per sub-question fused by RRF};
\node[opt, right=of ret] (rr) {\textbf{rerank}\\(if \code{use_reranker})\\cross-encoder, all 20};
\node[proc, right=of rr] (cut) {\textbf{narrow}\\first 5};
\node[opt, right=of cut] (cg) {\textbf{guard\_context}\\(if \code{check_retrieved_context})\\drop directive passages};
\node[rec, below=of dec] (r1) {\code{trace.sub_questions}};
\node[rec, below=of ret] (r2) {\code{trace.retrieved_chunk_ids}\\scored by recall@k, hit@k};
\node[rec, below=of rr] (r3) {\code{trace.rerank}\\movement statistics};
\node[rec, below=of cut] (r4) {\code{trace.ranked_chunk_ids}\\scored by MRR; first 5 by recall@context};
\node[rec, below=of cg] (r5) {\code{trace.guard_context}};
\draw[arr] (dec) -- (ret); \draw[arr] (ret) -- (rr); \draw[arr] (rr) -- (cut); \draw[arr] (cut) -- (cg);
\draw[dashed] (dec) -- (r1); \draw[dashed] (ret) -- (r2); \draw[dashed] (rr) -- (r3); \draw[dashed] (cut) -- (r4); \draw[dashed] (cg) -- (r5);
\end{tikzpicture}}
\caption{The retrieval stage. Dashed boxes are optional steps; yellow boxes are what each step writes into the trace and which evaluation metric reads it.}
\label{fig:retrieve-stage}
\end{figure}

Two design points are visible in the code comments. The candidate ids are recorded \emph{before} narrowing, because \enquote{the gap between that and what reaches the prompt is the thing reranking exists to close}. And the context guard runs \emph{after} narrowing, because \enquote{only what reaches the prompt can influence the model} \ev{src/vidyarag/pipeline.py:115-135}.

\paragraph{Why 20 candidates and 5 passages.} \code{config/default.yaml} explains that a reranker \enquote{can only add value when given a candidate pool wider than the final selection}, and calls a 4$\times$ pool \enquote{the latency/quality sweet spot on 2 vCPU} \ev{config/default.yaml:56-65}. The sweet-spot claim is not backed by a committed sweep; no alternative pool size was measured.

\section{Dense retrieval}
\label{sec:dense}

\subsection{Code}

\begin{lstlisting}[language=Python]
def retrieve_dense(client, query, *, collection, embedding_model, limit=20, book_slug=None):
    vector = embed_texts([query], embedding_model)[0]

    query_filter = None
    if book_slug is not None:
        query_filter = models.Filter(must=[models.FieldCondition(
            key="book_slug", match=models.MatchValue(value=book_slug))])

    response = client.query_points(collection_name=collection, query=vector,
                                   using=DENSE_VECTOR, limit=limit,
                                   query_filter=query_filter, with_payload=True)
    return [RetrievedChunk.from_payload(point.payload or {}, float(point.score))
            for point in response.points]
\end{lstlisting}
\vspace{-4pt}{\raggedleft\ev{src/vidyarag/retrieve/dense.py:81-124} (docstring removed)\par}

The module calls itself \enquote{the frozen baseline}: \enquote{Deliberately the simplest thing that works\ldots{} Every enhancement in Phase 4 is measured as a delta against this, so it must stay boring} \ev{src/vidyarag/retrieve/dense.py:1-10}. It returns the whole payload rather than just text, because citations are assembled from that metadata two layers later.

\subsection{The mathematics}

The question $q$ and every passage $d$ are mapped by the same encoder $E$ to 768-dimensional vectors, and passages are ranked by cosine similarity:
\begin{equation}
\operatorname{sim}(q,d) \;=\; \cos\theta \;=\; \frac{E(q)\cdot E(d)}{\lVert E(q)\rVert\,\lVert E(d)\rVert} \;=\; \frac{\sum_{i=1}^{768} E(q)_i\,E(d)_i}{\sqrt{\sum_i E(q)_i^2}\;\sqrt{\sum_i E(d)_i^2}}.
\end{equation}
The collection is created with \code{Distance.COSINE} \ev{src/vidyarag/store/collection.py:87-95}, and the returned \code{score} is this similarity. The top \code{limit} passages by score are returned.

\begin{background}[bi-encoders]
A model like BGE is a \emph{bi-encoder}: it encodes the question and each passage separately. Passage vectors can therefore be computed once, at ingestion, and a query costs one encoder pass plus a vector search. The price is that the model never sees the question and a passage together, so it judges topical similarity well and fine-grained relevance less well. Chapter~\ref{ch:concepts} explains embeddings and cosine similarity from first principles.
\end{background}

\subsection{Details worth knowing}

\begin{itemize}
  \item \textbf{Queries and passages are embedded by the same call.} \code{retrieve_dense} uses \code{embed_texts}, the ingestion function, so no query-specific instruction prefix is added. During this analysis, adding BGE's recommended prefix left the top three results for \enquote{what is tissue?} unchanged \tagM.
  \item \textbf{The model must match the index.} Its docstring warns that a different model \enquote{produces vectors in an unrelated space, and the search silently returns nonsense rather than failing.} Only a different vector \emph{width} is caught, at ingestion.
  \item \textbf{The book filter exists but the pipeline never uses it.} \code{book_slug} is passed only by the \code{/v1/search} endpoint. The \code{/v1/query} request model accepts a \code{book_slug} field too, but the handler ignores it (Chapter~\ref{ch:interfaces}).
  \item \textbf{\code{RetrievedChunk.from_payload} is tolerant}: missing text fields become empty strings, missing numbers become 0 and missing optional fields become \code{None} \ev{src/vidyarag/retrieve/dense.py:48-78}.
\end{itemize}

\section{Cross-encoder reranking}
\label{sec:rerank}

\subsection{Why it was added}

The reranker's docstring states the case from the baseline run's numbers: recall@$k$ was 0.967 but recall@context only 0.880, so \enquote{the gold passage is in the retrieved pool for almost every question, and is then ranked out of the top 5 before generation for roughly one question in eight. Reranking does not need to find anything new. It needs to reorder what dense retrieval already found} \ev{src/vidyarag/retrieve/rerank.py:8-12}.

\begin{background}[cross-encoders]
A cross-encoder reads the question and one passage \emph{together}, as one input such as \code{[CLS] question [SEP] passage [SEP]}, and outputs a single relevance score. Attention across both texts lets it notice whether the passage actually answers the question, not just whether it is on the same topic. The cost is one full model pass per (question, passage) pair at query time, which is why it is applied only to a short candidate list. The MS MARCO models output unnormalised scores that can be negative; only their order matters.
\end{background}

\subsection{Code}

\begin{lstlisting}[language=Python]
@lru_cache(maxsize=2)
def get_reranker(model_name: str) -> TextCrossEncoder:
    from fastembed.rerank.cross_encoder import TextCrossEncoder
    return TextCrossEncoder(model_name=model_name)

def rerank(query, chunks, *, model_name, top_k=None):
    if not chunks:
        return []
    scores = list(get_reranker(model_name).rerank(query, [c.text for c in chunks]))
    rescored = [dataclasses.replace(chunk, score=float(score), prior_score=chunk.score)
                for chunk, score in zip(chunks, scores, strict=True)]
    rescored.sort(key=lambda c: c.score, reverse=True)
    return rescored if top_k is None else rescored[:top_k]
\end{lstlisting}
\vspace{-4pt}{\raggedleft\ev{src/vidyarag/retrieve/rerank.py:31-74} (docstrings removed)\par}

The pipeline reranks all 20 candidates and narrows to 5 afterwards (it never passes \code{top_k}). Each reranked passage keeps its dense score in \code{prior_score}, \enquote{so a rank change can be attributed rather than guessed at.}

\subsection{Measuring what the reranker did}

\begin{lstlisting}[language=Python]
def rank_movement(before, after) -> dict[str, Any]:
    before_rank = {chunk.chunk_id: index for index, chunk in enumerate(before)}
    moved = sum(1 for index, chunk in enumerate(after)
                if before_rank.get(chunk.chunk_id, index) != index)
    promoted_into_top = [chunk.chunk_id for index, chunk in enumerate(after[:5])
                         if before_rank.get(chunk.chunk_id, 0) >= 5]
    return {"reordered": moved,
            "promoted_into_context": len(promoted_into_top),
            "top1_changed": bool(before and after and before[0].chunk_id != after[0].chunk_id)}
\end{lstlisting}
\vspace{-4pt}{\raggedleft\ev{src/vidyarag/retrieve/rerank.py:77-97}\par}

The docstring states the principle: \enquote{A reranker that improves a metric while never changing the top 5 has not earned the credit, and this is what would reveal that.} The size of the context window is written into this function as the literal \code{5} rather than read from \code{top_k_context}, so the statistic would silently describe the wrong window if the context size were ever changed.

\begin{verified}[what reranking did on the 58-question gold set, from the committed run files]
\begin{center}
\begin{tabular}{@{}lccc@{}}
\toprule
& \code{baseline} & \code{rerank} & Change\\
\midrule
Recall@$k$ (pool of 20) & 0.967 & 0.967 & $\pm$0.000\\
Recall@context (first 5) & 0.880 & 0.913 & +0.033\\
MRR & 0.770 & 0.830 & +0.060\\
Recall@context, factual questions & 0.893 & 0.964 & +0.071\\
Recall@context, multi-hop questions & 0.861 & 0.833 & $-$0.028\\
\bottomrule
\end{tabular}
\end{center}
Recall@$k$ did not move, which is exactly what the \code{rerank} profile's comment says must happen if the reranker is the only change. Per query, the reranker reordered 17.7 of the 20 candidates on average, changed the top passage for 32 of 58 questions (55\%) and promoted at least one passage from outside the first five into the context for 55 of 58 (95\%). Sources: \code{eval/results/baseline\_\_20260819T175415Z.json}, \code{eval/results/rerank\_\_20260820T121447Z.json}; the per-type and movement figures were computed from those files during this analysis and match \code{docs/EVALUATION.md}.
\end{verified}

\paragraph{The cost.} Mean latency rose from 1,091\,ms to 6,856\,ms. The median rerank-profile latency was 3,100\,ms; the mean is pulled up by the first question of the run, which took 44,109\,ms while models loaded \tagM. In a steady-state benchmark on the development machine (Intel Core i5-12450H), reranking 20 passages took a median of 2,142\,ms and 5 passages 444\,ms; in-pipeline measurements at other times during the analysis ranged up to 10--15 seconds for 20 passages \tagM. \code{docs/DESIGN.md} documents about 130\,ms per batch on CPU. Chapter~\ref{ch:performance} reconciles these.

\paragraph{The repository's explanation for the multi-hop drop.} A cross-encoder \enquote{scores each passage independently against the query and has no notion of what else was selected, so it promotes whatever most resembles the question --- typically several near-duplicates of the strongest match. A factual question has one right passage and that is exactly correct. A multi-hop question needs two complementary passages, and the duplicates crowd the second one out} \ev{src/vidyarag/retrieve/decompose.py:3-11}. This mechanism is the stated hypothesis that motivated decomposition; the repository does not isolate it experimentally.

\section{Query decomposition with reciprocal rank fusion}
\label{sec:decompose}

\subsection{Idea}

If a multi-hop question needs two passages about different topics, ask two narrower questions, retrieve for each, and fuse the lists, so that each hop gets its own chance at the passage it needs \ev{src/vidyarag/retrieve/decompose.py:12-15}.

\subsection{Splitting the question}

\begin{lstlisting}[language=Python]
class Decomposition(BaseModel):
    is_multi_hop: bool = Field(description="True only if answering requires facts from two or more distinct topics.")
    sub_questions: list[str] = Field(default_factory=list, description="Two or three self-contained questions. Empty if not multi-hop.")

def decompose(llm, question, *, model) -> Decomposition:
    try:
        response = llm.models.generate_content(
            model=model,
            contents=DECOMPOSE_PROMPT.format(question=question),
            config={"temperature": 0.0,
                    "response_mime_type": "application/json",
                    "response_schema": Decomposition})
    except Exception:  # noqa: BLE001 - degrade to baseline, never fail the query
        return Decomposition(is_multi_hop=False)
    ...  # use response.parsed, else parse response.text; any failure -> not multi-hop

def _clean(decomposition):
    subs = [s.strip() for s in decomposition.sub_questions if s and s.strip()]
    subs = subs[:MAX_SUB_QUESTIONS]            # at most 3
    if len(subs) < 2:
        return Decomposition(is_multi_hop=False)   # one sub-question is a paraphrase
    return Decomposition(is_multi_hop=True, sub_questions=subs)
\end{lstlisting}
\vspace{-4pt}{\raggedleft\ev{src/vidyarag/retrieve/decompose.py:46-121} (abridged)\par}

The prompt asks whether answering needs facts from \enquote{two or more DISTINCT topics that would live in different sections of a textbook}, requires each sub-question to be self-contained, and warns that \enquote{Splitting an atomic question produces near-duplicates that retrieve the same passage repeatedly and crowd out everything else, which is worse than not splitting at all} \ev{src/vidyarag/retrieve/decompose.py:58-75}. Gemini is asked for JSON matching the \code{Decomposition} schema. Every failure --- an exception, an empty response, invalid JSON, fewer than two sub-questions --- falls back to \enquote{not multi-hop}, which means plain dense retrieval: \enquote{the safe direction}.

\subsection{Fusing ranked lists}

Each sub-question retrieves its own top 20 (\enquote{the full \code{limit}, not a share of it}, so that a drop in recall cannot be blamed on a smaller budget), and the lists are fused by \emph{reciprocal rank fusion} \ev{src/vidyarag/retrieve/decompose.py:144-169}:
\begin{equation}
\operatorname{RRF}(d) \;=\; \sum_{i\,:\,d\in L_i} \frac{1}{k + r_i(d)}, \qquad k = 60,
\end{equation}
where $L_i$ is the ranked list for sub-question $i$ and $r_i(d)$ is the 1-based rank of passage $d$ in that list. Passages are sorted by fused score, and the pipeline keeps the first 20.

\begin{lstlisting}[language=Python]
def reciprocal_rank_fusion(ranked_lists, *, k=RRF_K):
    scores: dict[str, float] = {}
    seen: dict[str, RetrievedChunk] = {}
    for ranked in ranked_lists:
        for rank, chunk in enumerate(ranked, start=1):
            scores[chunk.chunk_id] = scores.get(chunk.chunk_id, 0.0) + 1.0 / (k + rank)
            seen.setdefault(chunk.chunk_id, chunk)
    ordered = sorted(scores.items(), key=lambda item: item[1], reverse=True)
    return [seen[chunk_id] for chunk_id, _ in ordered]
\end{lstlisting}
\vspace{-4pt}{\raggedleft\ev{src/vidyarag/retrieve/decompose.py:124-141} (docstring removed)\par}

\begin{example}[agreement beats one strong hit --- illustrative ranks]
Passage A is ranked 1st by sub-question 1 and absent from list 2: $\operatorname{RRF}(A) = \tfrac{1}{61} = 0.01639$. Passage B is ranked 5th by sub-question 1 and 4th by sub-question 2: $\operatorname{RRF}(B) = \tfrac{1}{65} + \tfrac{1}{64} = 0.01538 + 0.01563 = 0.03101$. B outranks A by almost a factor of two. A test encodes exactly this property: \code{test_agreement_across_sub_questions_outranks_one_strong_hit}.
\end{example}

The constant $k=60$ is \enquote{the value from the original paper}, deliberately left untuned \enquote{because tuning it on the gold set would be fitting the evaluation} \ev{src/vidyarag/retrieve/decompose.py:33-39}. RRF uses ranks only because dense scores from two different sub-queries \enquote{are not on a comparable scale} \ev{src/vidyarag/retrieve/decompose.py:17-21}.

\begin{inferred}[two properties of this implementation]
The fused list keeps the \code{RetrievedChunk} object from the first list a passage appeared in, so its \code{score} is that sub-question's cosine similarity, not the fused score. Scores shown in traces or API responses for decomposed retrieval are therefore not the values that ordered the list. And ties are broken by first appearance, because Python's sort is stable and the score dictionary preserves insertion order.
\end{inferred}

\subsection{What the measurement showed}
\label{sec:decompose-result}

\begin{verified}[the \code{decompose} run against \code{baseline}, from the committed run files]
Decomposition split 26 of the 58 questions: 6 of 28 factual (21\%), 11 of 18 multi-hop (61\%) and 9 of 12 unanswerable (75\%). Overall recall@$k$ fell from 0.967 to 0.957, recall@context from 0.880 to 0.826 and MRR from 0.770 to 0.769. On multi-hop questions, pool recall rose to 1.000, meaning every gold passage was among the 20 candidates, while recall@context fell from 0.861 to 0.778. On the 11 multi-hop questions that were actually split, recall@context was 0.727, against 0.864 for the same questions under \code{baseline}. Split rates and per-type figures were computed from \code{eval/results/decompose\_\_20260824T065521Z.json} during this analysis and agree with \code{docs/EVALUATION.md}.
\end{verified}

The repository's diagnosis is that RRF rewards consensus: a passage retrieved by both sub-questions is usually the less specific one, while a passage that only one hop needs appears in only one list and is pushed out of the first five. \code{docs/EVALUATION.md} names the fix it did not build --- reserving prompt slots for each sub-question --- and records the profile as rejected. The code stays in the repository so that the negative result remains reproducible (\code{.env.example} describes \code{decompose} as \enquote{Measured and rejected}).

The stored sub-questions also show the failure the prompt warns about. For \code{fact-003}, a single-lookup factual question, the model produced two near-duplicate sub-questions \tagM.

\begin{interview}
A rejected technique is one of the strongest things you can discuss. The complete story: the reranker helped factual questions and hurt multi-hop ones; the stated hypothesis was a diversity failure; decomposition was the targeted fix; it put every multi-hop gold passage into the candidate pool (pool recall 1.000) but then lost them in the final ordering (recall@context 0.778), because rank fusion favours passages that several sub-questions agree on. It was rejected on the numbers, and the obvious next step --- a guaranteed slot per sub-question --- is known and was not built.
\end{interview}

\section{Hybrid search: reserved, not built}

Sparse (BM25) retrieval fused with dense retrieval is the most common RAG upgrade, and the configuration has keys for it: \code{use_hybrid} and \code{sparse_model: Qdrant/bm25}. It was never built. The recorded reason is that first-stage recall@$k$ was already 0.967, \enquote{leaving hybrid retrieval almost nothing to recover}, and that adding it would require re-indexing the whole corpus with sparse vectors \ev{src/vidyarag/store/collection.py:29-38}. Rather than leave a switch that does nothing, the configuration refuses \code{use_hybrid: true} at load time \ev{src/vidyarag/settings.py:164-181}.

\begin{inferred}
The recall argument covers the gold set's questions, whose gold passages dense retrieval already finds. It does not cover queries that depend on exact terms --- rare technical words, abbreviations, names --- where sparse retrieval traditionally helps, and the gold set was not designed to measure that. \code{fact-002} is such a case: its answer is a single term (\enquote{chemotaxis}), and its gold passage was not among the 20 dense candidates at all.
\end{inferred}

\section{How evaluation reads the retrieval trace}

Three lists recorded during retrieval feed three different metrics (Chapter~\ref{ch:evaluation}):

\begin{center}
\begin{tabularx}{\textwidth}{@{}>{\raggedright\arraybackslash}p{3.8cm}>{\raggedright\arraybackslash}p{4.0cm}L@{}}
\toprule
\textbf{Trace list} & \textbf{Metrics} & \textbf{Why that list}\\
\midrule
\code{retrieved_chunk_ids} & recall@$k$, hit@$k$ & Did first-stage retrieval find the gold passage at all?\\
\code{ranked_chunk_ids} & MRR & Where did the gold passage end up in the final ordering?\\
first 5 of the ranking & recall@context & Did the gold passage reach the prompt?\\
\bottomrule
\end{tabularx}
\end{center}

This separation is itself a bug fix. \code{evaluation/retrieval.py} records that the first rerank ablation scored every metric on the pool list, so context recall and MRR came out \emph{bit-identical} to the baseline while context precision moved six points: \enquote{the instrument read the wrong list} (Chapter~\ref{ch:history}).

\section{Retrieval findings and limits}

\begin{itemize}
  \item \textbf{A hard-coded window.} \code{rank_movement} uses the literal 5 instead of \code{top_k_context} \ev{src/vidyarag/retrieve/rerank.py:90-91}. \tagV
  \item \textbf{Fused scores are not reported.} Decomposed results carry the first-seen dense score. \tagI
  \item \textbf{Retries overwrite the trace.} In the corrective loop each retry calls \code{retrieve} again with the same trace object, so \code{retrieved_chunk_ids} and \code{ranked_chunk_ids} describe only the last attempt \ev{src/vidyarag/pipeline.py:213-216}. \tagV
  \item \textbf{Decomposition spends quota invisibly.} It uses the generation model, so it draws on the same free-tier limits, and its token usage is never added to the trace (Section~\ref{sec:token-bug}). \tagV
  \item \textbf{Local reranking is slow on the development CPU} (about 10 seconds for 20 passages, measured). \tagM
  \item \textbf{Short definitional questions} favour section text over glossary definitions (Section~\ref{sec:chunking}). \tagM
\end{itemize}

\begin{recommended}
Read the window size from configuration in \code{rank_movement}. Carry the fused RRF score on decomposed results. Keep per-attempt candidate lists in the trace. If multi-hop questions matter, try decomposition with one guaranteed context slot per sub-question, or maximal marginal relevance over the reranked pool, measured against the frozen baseline in the usual way. Evaluate a hybrid variant on a query set that contains exact-term questions before concluding that sparse retrieval is unnecessary.
\end{recommended}
